%======================================================================
% 3. UNCERTAINTY MODEL AND SURROGATE CONSTRUCTION
%======================================================================
\section{Uncertainty model and surrogate construction}
\label{sec:surrogate}

\subsection{Uncertain variables}
\label{sec:uncertainty}

Four inputs are uncertain (Table~\ref{tab:uncert}). Two describe the
entry-interface state: the flight-path angle
$\gamma_0\sim\mathcal{N}(-2^\circ,\,0.5^\circ)$ and the entry speed
$V_0\sim\mathcal{N}(7830,\,50)$~m/s, the nominal values being those of the
underlying optimization case. The magnitudes --- half a degree on
flight-path angle, 0.6\% on speed --- are representative of
entry-interface targeting accuracy for guided lunar-return entries, but
they are assumptions, not measurements: Section~\ref{sec:sens} therefore
brackets $\sigma_{\gamma_0}$ over a factor of three, and every comparison
conclusion survives that range. Two represent aerodynamic model form:
multiplicative factors $k_{CD},k_{CL}\sim\mathcal{N}(1,\,0.05)$ applied to
the drag and lift interpolators. The factors perturb force magnitudes at
fixed trim (the trim angle comes from the unscaled moment), which isolates
force-level model error from trim prediction. All four are treated as
independent; sampling is by Latin hypercube through the inverse CDF.
Independence of the two scale factors is the weakest part of this
specification --- the fitted Mach-wise scale factors move together across
Mach --- and Section~\ref{sec:sens} therefore probes every headline
statistic against correlated scales, wider and narrower scale bands, and
alternative entry-state dispersions.

The 5\% band on $k_{CD}$ and $k_{CL}$ is sized from the validation
residuals of Section~\ref{sec:validation} rather than assumed.
Table~\ref{tab:calib} shows the calibration audit explicitly: for each
force coefficient at the two validation Mach numbers, the median scale
factor that best aligns the model with the digitized curves, its distance
from unity in units of the assumed $\sigma_k=0.05$, and the residual shape
error that remains after that scale is removed. Two features size the
model. First, the post-scale residuals are small: at most 0.013
(0.9--2.2\% of the coefficient magnitude) in three of the four
coefficient/Mach combinations, with $C_D$ at $M=3$ the outlier at 0.060
(4.4\%), the documented tangent-cone beyond-detachment case --- so a
multiplicative band of a few percent covers the residual error. Second,
the fitted scales are themselves Mach-structured: near 0.95 at $M=10$
($\approx1\sigma$) but 1.11 and 1.19 for $C_D$ and $C_L$ at $M=3$ (2.2
and 3.7$\sigma$). A single trajectory-constant Gaussian multiplier cannot
represent that structure exactly, and we do not claim it does: the 5\%
standard deviation covers the residual error comfortably and treats the
$M=3$ calibration bias --- which a trajectory-constant factor would
partially average over the entry Mach range --- as a tail event rather
than a central one. What protects the conclusions is that they do not
depend on this choice: Section~\ref{sec:sens} varies $\sigma_k$ by
$\pm40\%$, correlates the two factors, and brackets $\sigma_{\gamma_0}$ by
a factor of three, with every headline statistic moving within a few
percentage points. The uncertainty model is thus \emph{sized from}
validation evidence in a traceable sense --- an exercise in
calibration-informed uncertainty in the spirit of Kennedy and O'Hagan
\cite{kennedy2001} --- not calibrated in the strict sense of a fitted
posterior.

\begin{table}[htbp]
\centering
\caption{Calibration audit of the aerodynamic scale factors: median
  fitted scale $k$ against the digitized Apollo coefficient curves, the
  corresponding distance from unity in units of $\sigma_k=0.05$, and the
  residual shape error after scale removal (max over angle of attack;
  percentage of the typical coefficient magnitude over the sweep).}
\label{tab:calib}
\begin{tabular}{lccccc}
\toprule
coefficient & $M$ & fitted $k$ & $(k-1)/\sigma_k$ &
post-scale residual & (\% of $|c|$) \\
\midrule
$C_D$ & 3  & 1.11 & $+2.2\,\sigma$ & 0.060$^{\dagger}$ & 4.4 \\
$C_L$ & 3  & 1.19 & $+3.7\,\sigma$ & 0.009 & 3.1 \\
$C_D$ & 10 & 0.95 & $-1.1\,\sigma$ & 0.013 & 0.9 \\
$C_L$ & 10 & 0.94 & $-1.1\,\sigma$ & 0.007 & 2.2 \\
\bottomrule
\end{tabular}

\smallskip
{\footnotesize $^{\dagger}$ documented tangent-cone beyond-detachment
method difference (Appendix~\ref{app:deviations}), not random error.}
\end{table}

\begin{table}[htbp]
\centering
\caption{Uncertain inputs. Entry-state dispersion plus aerodynamic
model-form scales; all Gaussian and independent.}
\label{tab:uncert}
\begin{tabular}{lccl}
\toprule
input & distribution & role \\
\midrule
$\gamma_0$ & $\mathcal{N}(-2^\circ,\ 0.5^\circ)$ & entry state \\
$V_0$      & $\mathcal{N}(7830,\ 50)\ \si{\metre\per\second}$ & entry state \\
$k_{CD}$   & $\mathcal{N}(1,\ 0.05)$ & drag model form \\
$k_{CL}$   & $\mathcal{N}(1,\ 0.05)$ & lift model form \\
\bottomrule
\end{tabular}
\end{table}

\subsection{Design of experiments}
\label{sec:doe}

The simulator was run 2{,}800 times over the seven inputs (three design,
four uncertain), sampled by Latin hypercube. The set grew in three stages: a
1{,}500-point base design spanning the full space, a 300-point enrichment of
the lofted corner (shallow $\gamma_0$ with small $r_\theta$), where early
surrogate errors concentrated, and a final 1{,}000-point enrichment of the
fast-descent family defined below. A run costs roughly 19~s of wall time.
Of the designs sampled, 44.4\% satisfy the constraints of
Section~\ref{sec:formulation} at nominal conditions; the remainder populate
the infeasible region and are retained with sentinel values so the surrogate
sees the whole space.

Two data-quality findings shaped the final set. First, trajectories that do
not reach the $M=3$ terminal event within the integration horizon are
silently truncated, and their recorded range and heat load then describe an
arbitrary intermediate state. An initial horizon of 7{,}200~s capped 202
lofted runs that in fact capture between 7.2 and 13.5~ks; the horizon was
raised to $t_{\max}=\SI{14400}{\s}$ (largest observed capture 13{,}474~s)
and the affected rows were re-run and overwritten, so every reported
objective is a terminal value. Second, the capture time is strongly bimodal
in log space (an empty gap between 2{,}241 and 2{,}297~s in the 1{,}800-run
set), splitting entries into a \emph{fast} family (direct descent, capture
in 0.4--1.9~ks) and a \emph{slow}, lofted family (4.1--13.5~ks), with the
boundary at 2{,}250~s. Descriptor quantities recorded for every run (capture
flag, final time, peak altitude after 60~km, minimum flight-path angle)
support the split: 1{,}795 of the first 1{,}800 runs capture, so capture
itself is near-certain and the operative structure is fast versus slow.

\subsection{Gaussian-process surrogates}
\label{sec:gp}

Each output ($Q_s$, $sg$, $q_{\mathrm{stag}}^{\max}$,
$q_{\mathrm{shldr}}^{\max}$, $n_{\max}$) is modelled in log space by a
Gaussian process --- the standard regression of the computer-experiments
tradition \cite{sacks1989} --- with a Mat\'ern-5/2 kernel, automatic
relevance determination (ARD) on all seven inputs, and a white-noise term,
fitted with scikit-learn \cite{rasmussen2006gp,pedregosa2011sklearn}. Range is not
modelled by a single process. A linear logistic classifier first assigns the
trajectory family from the seven inputs (out-of-fold AUC 0.9994, accuracy
0.989), and a dedicated GP is fitted to $\log sg$ within each family; the
mixture prediction combines the two through the family probabilities. A
companion capture classifier (AUC 0.99) guards the rare non-capturing
corner. The two geometry constraints depend on the shape only, but their
exact evaluation requires rebuilding the aerodynamic database per candidate;
inside the optimizers they are therefore evaluated by lightweight models
trained on the DOE --- a gradient-boosted trim-existence classifier (98.5\%
out-of-fold; in the DOE, absent trim and static instability coincide
exactly) and a nose-margin regressor (98.8\% sign accuracy within trim).

The slow family deserves comment, because it is where surrogate
construction usually stops being routine. Within that family $\log sg$ is
not learnable at any practical design density: Gaussian-process,
gradient-boosted and random-forest regressions reach log-space $R^2$ of
0.66, 0.41 and 0.45 on held-out data, single-input correlations are near
zero, and the residual noise floor is 0.11 in $\log sg$. Perturbation probes
show why. Around a slow design, the response to $\gamma_0$ is locally linear
and smooth, with $\mathrm{d}(sg)/\mathrm{d}\gamma_0\approx
-2\times10^{5}$~km/degree --- roughly forty times steeper than the
fast-family sensitivity ($+5\times10^{3}$~km/degree) and of opposite sign.
A thousandth of a degree moves the predicted range by about 200~km. The
response is steep, not chaotic; slow-family range dispersion is physical
sensitivity that no practicable sampling density resolves, and it is carried
as an explicit dispersion term rather than averaged away. This is itself a
result in the spirit of robust design: the quantity that the deterministic
problem treats as a deterministic objective is, for lofted entries, a
dispersion-dominated one.

\subsection{Predictive-accuracy gate}
\label{sec:gate}

Before any optimization, the surrogates were passed through a five-fold
out-of-fold gate with per-output error bands fixed in advance
(Table~\ref{tab:gate}). Two outputs pass cleanly and two sit at their bands
(within 0.3 and 0.01 percentage points of the respective limits); fast-family
range passes at 3.8\% of range. Slow-family range is reported as what it is:
irreducible at practical sampling density, for the reason above. Fitted
white-noise levels of 1--3\% in log space are consistent with the gate
errors. The ARD length-scales add a physical footnote: for both peak-heating
outputs, the length-scales on $k_{CD}$ and $k_{CL}$ exceed the normalized
input range by more than an order of magnitude --- model-form uncertainty
concentrates in the integrated objectives ($Q_s$, $sg$) and leaves the peaks
nearly untouched.

\begin{table}[htbp]
\centering
\caption{Out-of-fold predictive accuracy of the surrogate models
  (5-fold CV, $n=2800$; RMSE as \% of the DOE range; log-space fits).
  sg uses the family-mixture construction.}
\label{tab:gate}
\begin{tabular}{lccc}
\toprule
output & $R^2$ & RMSE (\% range) & verdict \\
\midrule
$q_{\mathrm{stag}}$ peak & 0.969 & 2.29 & pass \\
$n_{\max}$              & 0.961 & 2.23 & pass \\
$q_{\mathrm{shldr}}$ peak & 0.952 & 3.01 & at 3\% band \\
$Q_s$                   & 0.983 & 2.33 & at 2\% band \\
$sg$ (fast family)      & 0.932 & 3.83 & pass \\
$sg$ (slow family)      & 0.562 & $\sim$15 & irreducible \\
\bottomrule
\end{tabular}
\end{table}
